%! The Statistical Cycle and Types of Data — Unit 1, Lesson 1.1 — Group Activity (Answer Key)
\documentclass[10pt]{article}
\usepackage{saar-article}
\usepackage{saar-key}   % pulls in -boxes; do NOT also load -boxes
\newcommand{\TallMath}[1]{$\displaystyle #1\rule[-1.4em]{0pt}{3.2em}$}

\begin{document}

\pageheader{Unit 1, Lesson 1.1}{Group Activity --- Answer Key}
% NAMESTRIP: no \namedateperiod here — mirrors the blank. Teacher-only prose goes
% in the lesson plan as \begin{teachernote}[Group Activity], never in a key.

\vspace{0.15in}

\begin{headlinebox}{frost}
\small
\textbf{Run the Cycle --- The Riverbend Breakfast Cart.} The cafeteria manager put
a breakfast cart in the main hallway for one week and wants to know whether to
keep it. Every answer needs a \emph{reason}, and every conclusion has to name the
group it applies to.
\end{headlinebox}

\vspace{0.1in}

\boxguard[30]
\begin{tcolorbox}[colback=white, colframe=black!40, breakable,
    title=\textbf{Tier R --- Remediate},
    fonttitle=\bfseries, arc=2mm, left=3mm, right=3mm, top=2mm, bottom=2mm]
\small
On Friday the manager asked $40$ students leaving the cart, ``What did you buy at
the cart this morning?''

\begin{enumerate}[leftmargin=1.6em, itemsep=6pt, topsep=4pt]

  \item Match each part of the manager's study to its stage. Write $1$, $2$, $3$,
        or $4$.

        \begin{center}
        \renewcommand{\arraystretch}{1.4}
        \begin{tabularx}{0.95\linewidth}{@{} X c @{}}
        \toprule
        \textbf{What the manager did} & \textbf{Stage} \\
        \midrule
        Asked 40 students what they bought.            & \ans{2} \\
        Decided to ask what students buy at the cart.  & \ans{1} \\
        Told the principal what the results showed.    & \ans{4} \\
        Counted the answers into a table by item.      & \ans{3} \\
        \bottomrule
        \end{tabularx}
        \end{center}

  \item Complete the frequency table. The first row is done for you.

        \begin{center}
        \renewcommand{\arraystretch}{1.4}
        \begin{tabularx}{0.9\linewidth}{@{} l c X @{}}
        \toprule
        \textbf{What they bought} & \textbf{Count} &
        \textbf{Relative frequency (percent)} \\
        \midrule
        Bagel      & 14 & $14/40 = 0.35 = 35\%$ \\
        Fruit cup  & 10 & \ans{$10/40 = 0.25 = 25\%$} \\
        Muffin     & 12 & \ans{$12/40 = 0.30 = 30\%$} \\
        Nothing    &  4 & \ans{$4/40 = 0.10 = 10\%$} \\
        \bottomrule
        \end{tabularx}
        \end{center}

  \item What was the \textbf{variable of interest} --- what did the manager
        record for each student?

        Variable: \ans{what that student bought} \hfill Type: \ans{categorical}

  \item Which item was bought by the most students? \ans{the bagel (35\%)}

  \item Write one sentence saying what the study found. Name the group and the
        percent.
        \ansline{Of the 40 students asked at the cart, 35\% bought a bagel that morning.}
\end{enumerate}
\end{tcolorbox}

\vspace{0.1in}

\boxguard
\begin{tcolorbox}[colback=white, colframe=black!40, breakable,
    title=\textbf{Tier A --- Approaching Proficiency},
    fonttitle=\bfseries, arc=2mm, left=3mm, right=3mm, top=2mm, bottom=2mm]
\small
\begin{enumerate}[leftmargin=1.6em, itemsep=5pt, topsep=3pt]

  \item The manager drafts two questions for a second study.

        \begin{itemize}[leftmargin=1.6em, itemsep=2pt, topsep=2pt]
          \item[\textbf{A.}] ``Is the breakfast cart good?''
          \item[\textbf{B.}] ``How many mornings per week do Riverbend students
                buy breakfast at school?''
        \end{itemize}

        \textbf{(a)} Which one is a statistical question? \ans{B}

        \textbf{(b)} Name one test the other question fails, and repair it.
        \ansline{A never says what to record. Repair: what percent of students buy breakfast here?}

        \textbf{(c)} For question \textbf{B}, what is the variable of interest,
        and what type is it?

        Variable: \ans{mornings per week} \hfill Type: \ans{quantitative}

  \item The manager surveys $80$ Riverbend students. Of those, $24$ say they
        would buy breakfast at school every morning if the cart stayed.

        \textbf{(a)} What percent of the students surveyed is that?

        \begin{work}
          \dfrac{24}{80} &= 0.30 \\
                         &= 30\%
        \end{work}

        \textbf{(b)} Riverbend has $900$ students. Predict how many would buy
        breakfast every morning.

        \begin{work}
          0.30 \times 900 &= 270 \text{ students}
        \end{work}

        \textbf{(c)} Write the conclusion as one sentence the manager could hand
        to the principal. Name the group, the number, and what was measured.
        \ansline{About 270 of Riverbend's 900 students say they would buy breakfast daily.}

  \item Each report below fails at one stage of the cycle. Write the stage number
        and say in a few words what went wrong.

        \begin{center}
        \renewcommand{\arraystretch}{1.5}
        \begin{tabularx}{0.97\linewidth}{@{} X c p{4.2cm} @{}}
        \toprule
        \textbf{Report} & \textbf{Stage} & \textbf{What went wrong} \\
        \midrule
        Only students standing at the cart were asked, and the report says
        ``Riverbend students love the cart.''
          & \ans{2} & \ans{only cart users were asked, so the data cannot describe everyone} \\
        The manager wrote down all 40 answers as a paragraph and never counted
        them by item.
          & \ans{3} & \ans{never organized, so no pattern is visible} \\
        \bottomrule
        \end{tabularx}
        \end{center}
\end{enumerate}
\end{tcolorbox}

\vspace{0.1in}

\boxguard
\begin{tcolorbox}[colback=white, colframe=black!40, breakable,
    title=\textbf{Tier E --- Extension},
    fonttitle=\bfseries, arc=2mm, left=3mm, right=3mm, top=2mm, bottom=2mm]
\small
\begin{enumerate}[leftmargin=1.6em, itemsep=5pt, topsep=3pt]

  \item The school newspaper prints this:

        \begin{center}
        \textit{``$35\%$ of Riverbend students buy a bagel every morning.''}
        \end{center}

        The arithmetic behind the $35\%$ is correct. Explain \emph{two} separate
        things the sentence claims that the Tier R data cannot support.
        \ansline{The 35\% describes only the 40 students leaving the cart, not all of Riverbend.}
        \ansline{And they were asked about one morning, so nothing supports ``every morning.''}

  \item Your group gets one hour on the district agenda. Choose a question about
        students at your school that data could settle, then plan all four
        stages.

        \begin{center}
        \renewcommand{\arraystretch}{1.35}
        \begin{tabularx}{0.97\linewidth}{@{} p{3.3cm} X @{}}
        \toprule
        Stage 1 --- our question & \ans{Answers vary --- must name the group and what is recorded.} \\
        Variable of interest & \ans{Answers vary --- one characteristic per student.} \\
        Type of data & \ans{Must match the variable named directly above.} \\
        Stage 2 --- who we ask, and how many & \ans{Answers vary --- a stated group and a stated number.} \\
        Stage 3 --- how we organize it & \ans{A frequency table of counts and percents, or a mean.} \\
        Stage 4 --- the sentence we hope to write & \ans{Must name the group, the number, and the units.} \\
        \bottomrule
        \end{tabularx}
        \end{center}

  \item Two groups ask about the same students. Group 1 records \emph{how many
        mornings per week} each student buys breakfast. Group 2 records only
        whether each student \emph{ever} buys breakfast --- yes or no.

        Type for Group 1: \ans{quantitative} \hfill Type for Group 2: \ans{categorical}

        The two groups collected data from the same students on the same topic.
        Explain how the \emph{question} --- not the students --- decided the type.
        \ansline{Group 1 asked for a count of mornings, an amount you can average.}
        \ansline{Group 2 asked for a yes-or-no label, which can only be counted.}
\end{enumerate}
\end{tcolorbox}

\end{document}
